% GENERATED by CasDesignBasis/tools/design_report.py from design_basis.toml v2.3.12. Do not edit.
\begin{tabularx}{\textwidth}{@{}l r l l L@{}}
\toprule
Symbol & Value & Unit & Source & Reason \\
\midrule
$K_{c}$ & \num{450} & rpm/m & [HEIM] ConfigValues & LIC01 proportional gain (rpm of pump speed per metre over-height); [FD] P = 1 normalised (90 rpm per 0.20 m). Util.PID form $u = K_c\big(e + \frac{1}{T_i}\int e\,dt\big)$ [OGA] 2-3. \\
$T_{i}$ & \num{60} & s & [HEIM] ConfigValues & LIC01 integral time, chosen for robustness (worst-case phase margin in the design basis section on LIC01, [OGA] 7-7); [FD] gives 600 s, about 20 times the process time constant. \\
$T_{d}$ & \num{0} & s & [HEIM] ConfigValues & No derivative action (PI): the level reading is noisy and quantized. \\
$f_c$ & \num{0.05} & Hz & [HEIM], [PRO] 10.3.3--10.3.4 & Cut-off of the 2nd-order Butterworth low-pass (bilinear transform with pre-warping, SCL.Butterworth) on each level transmitter; ten times above the loop crossover, attenuates waves of 4--10 s period by about 19 dB. \\
$\theta$ & \num{2} & s & - & Dead time in the level loop: Modbus TCP poll, task cycle, OPC UA to the Framo PLC and the drive. The margins are also given for $2.5\theta$.\assumed{} \\
$T_{s,\text{O}_2}$ & \num{24} & s & [HEIM] ModbusRTUData & Time between two readings of the same Orbit 864: five sensors polled one at a time over the 750-652 (2--3 s each), the oxygen block every second round. The sensor lag $\tau_s$ acts as the anti-aliasing filter [PRO] 6.2.\assumed{} \\
\bottomrule
\end{tabularx}
